\chapter{Derivatives and Partial Derivatives}
\label{sec:derivatives}

Chapter~\ref{sec:regular-expressions} introduced the syntax and semantics
of regular expressions. The semantics provides a declarative definition of
membership: a word $w$ matches an expression $r$ exactly when
$w \in L(r)$. This chapter introduces three algorithms for deciding membership.

The first is a direct "naive" recursive matcher
that follows the structure of the semantic definition. The second uses
Brzozowski derivatives to transform an expression as input symbols are
consumed. The third uses Antimirov's partial derivatives, representing
the possible residual expressions as a collection rather than combining
them into a single expression. All three are implemented in Rust with the
common interface \texttt{fn(\&str, \&Regex) -> bool}.

The three approaches decide the same membership problem, but they differ
in how they represent and explore the remaining matching problem. The
derivative-based approaches introduced here form the basis of the parsers developed in
Chapters~\ref{sec:derivative-parser} and
\ref{sec:partial-derivative-parser}.

%

\section{Deciding Membership}
\label{sec:membership-overview}

The semantics of Section~\ref{sec:regex-semantics} defines membership
declaratively: a word $w$ matches a regular expression $r$ exactly when
$w \in L(r)$. A direct way to decide membership is therefore to interpret
these semantic definitions as recursive matching rules.

Let $w$ be represented as an indexable sequence of characters, and let $w[i:j]$
denote the substring beginning at position $i$ and ending immediately
before position $j$. A recursive procedure
$\operatorname{match}(w,i,j,r)$ can then determine whether
$w[i:j] \in L(r)$ by following the structure of $r$:


\begin{itemize}
  \item $\varepsilon$ matches $w[i:j]$ iff $i = j$.
  \item A literal $x$ matches $w[i:j]$ iff $j = i+1$ and the character at
        position $i$ is $x$.
  \item $r + s$ matches $w[i:j]$ iff $r$ matches $w[i:j]$ or $s$ does.
  \item $r \cdot s$ matches $w[i:j]$ iff there is some split point $k \in
        [i,j]$ such that $r$ matches $w[i:k]$ and $s$ matches $w[k:j]$.
  \item $r^{*}$ matches $w[i:j]$ iff $i = j$, or there is some $k \in
        (i,j]$ such that $r$ matches $w[i:k]$ and $r^{*}$ matches
        $w[k:j]$.
\end{itemize}


Membership of the whole word $w$ corresponds to $\operatorname{match}(w,
0, |w|, r)$. Each matching case is a direct operational restatement 
of the corresponding clause in the definition of $L(r)$, 
meaning the correctness follows directly from the regular-expression semantics.

It is also, in general, an expensive procedure. Concatenation and star
both search over every split point in the current range rather than
committing to one, and nothing remembers a range once it has already
been explored while resolving a different call. Take $r = a \cdot a
\cdot a$ against $w = \mathtt{aaa}$. Matching $w[0:3]$ against $r$ tries
every split point $k \in \{0,1,2,3\}$ for the outer concatenation, even
though only $k = 1$ can possibly succeed, since the left operand is a
single literal and can only match a range of length exactly $1$. The
right operand at that point, $w[1:3]$ against $a \cdot a$, repeats the
same search one level down. For a chain of $n$ literals against an input
of length $n$, the number of ranges examined this way grows
exponentially in $n$ rather than linearly, since each concatenation node
re-explores the full range handed to it regardless of what a sibling
call already tried.

\texttt{src/matchers/match\_naive.rs} implements this procedure directly
on a \texttt{Vec<char>} and a pair of \texttt{usize} bounds. Despite its
cost, it remains useful: because it transcribes $L(r)$'s definition
without any intervening transformation, it serves as an independent
reference that the more efficient derivative-based matchers introduced below can be
checked against. 

The derivative approach takes a different view of the same problem.
Rather than searching for a decomposition of $w$ while keeping $r$ fixed,
it \emph{transforms} $r$ as each character of $w$ is consumed. The
resulting expression describes precisely the language that can still
match the remaining suffix, allowing the matcher to continue with an
adjusted expression rather than reconsidering the original $r$ at every
step.

%

\section{Brzozowski Derivatives}
\label{sec:brzozowski-derivatives}

Brzozowski's derivative of $r$ with respect to a symbol $x$, written
$r \backslash x$, is the expression describing exactly what remains of
$r$'s language once a leading $x$ has been removed \citep{brzozowski1964}:

\begin{equation}\label{eq:der-23}
L(r \backslash x) \;=\; \{\, w \mid x w \in L(r) \,\}.
\end{equation}

An expression is \emph{nullable} if it accepts the empty word
$\varepsilon$. Let $n(r)$ denote the corresponding Boolean predicate, so
that $n(r)$ is true exactly when $\varepsilon \in L(r)$. Deciding whether a
word $x_1 x_2 \cdots x_n$ is in $L(r)$ then reduces to computing the chain
$r \backslash x_1 \backslash x_2 \backslash \cdots \backslash x_n$ and
checking whether the final expression is nullable. Both nullability and
the derivative itself can be defined by structural recursion directly on
the syntax of $r$, with no reference back to $L(r)$ needed once the base
cases are fixed. Equation~\eqref{eq:nullability} gives nullability and
Equation~\eqref{eq:derivative} the derivative:


\begin{equation}\label{eq:nullability}
\begin{aligned}
n(\emptyset) &= \mathrm{false} &&(\mathsf{N}_{\emptyset}) &
n(r + s) &= n(r) \lor n(s) &&(\mathsf{N}_{+}) \\
n(\varepsilon) &= \mathrm{true} &&(\mathsf{N}_{\varepsilon}) &
n(r \cdot s) &= n(r) \land n(s) &&(\mathsf{N}_{\mathrm{seq}}) \\
n(x) &= \mathrm{false} &&(\mathsf{N}_{\mathrm{lit}}) &
n(r^{*}) &= \mathrm{true} &&(\mathsf{N}_{*})
\end{aligned}
\end{equation}
\begin{equation}\label{eq:derivative}
\begin{aligned}
\emptyset \backslash x &= \emptyset &&(\mathsf{D}_{\emptyset}) \\
\varepsilon \backslash x &= \emptyset &&(\mathsf{D}_{\varepsilon}) \\
y \backslash x &=
  \begin{cases} \varepsilon & x = y \\ \emptyset & x \neq y \end{cases}
  &&(\mathsf{D}_{\mathrm{lit}}) \\
(r + s) \backslash x &= (r \backslash x) + (s \backslash x) &&(\mathsf{D}_{+}) \\
r^{*} \backslash x &= (r \backslash x) \cdot r^{*} &&(\mathsf{D}_{*}) \\
(r \cdot s) \backslash x &=
  \begin{cases}
    (r \backslash x) \cdot s \; + \; s \backslash x & n(r) = \mathrm{true} \\
    (r \backslash x) \cdot s & n(r) = \mathrm{false}
  \end{cases}
  &&\begin{array}{l}(\mathsf{D}_{\mathrm{seq}}^{\,n}) \\ (\mathsf{D}_{\mathrm{seq}})\end{array}
\end{aligned}
\end{equation}

Each rule carries a name in the right-hand column, and every derivation
in this chapter and in
Chapters~\ref{sec:derivative-parser} and~\ref{sec:partial-derivative-parser}
cites those names step by step. Two of them, $\mathsf{D}_{+}$ and
$\mathsf{D}_{\mathrm{seq}}$, are stated over metavariables $r$ and $s$, so
applying one to a concrete expression means first deciding which
subexpression each metavariable stands for. Every application below
states that binding explicitly before the rule is used, because for a
left-associated expression such as
$r_{01} = ((0+1)^{*} \cdot 0) \cdot 1$ the answer is not the one the
written form suggests at a glance: the outer concatenation binds
$r = (0+1)^{*} \cdot 0$ and $s = 1$, not $r = (0+1)^{*}$ and
$s = 0 \cdot 1$.

The only case that needs justification beyond direct case analysis on
$L(r \cdot s)$ is concatenation, the last line of
Equation~\eqref{eq:derivative}: if $r$ is nullable, a word starting with
$x$ can either continue matching $r$ (so the derivative goes into $r$
while $s$ is left untouched) or arise from $r$ already being satisfied by
$\varepsilon$, with $x$ belonging to $s$ instead. Hence the extra
$s \backslash x$ summand appears exactly when $n(r)$ holds. If $r$ is not
nullable, this second possibility cannot arise, since $\varepsilon \notin
L(r)$ rules out $r$ having already been satisfied before any input was
consumed.

%

\paragraph{Worked example: single-symbol matching.}
Consider $x^{*}$ against the word $\mathtt{xx}$. Writing
$r \xrightarrow{x} r \backslash x$ for one derivative step,

\emph{Step 1.} The expression is a star, so $\mathsf{D}_{*}$ applies with
body $r = x$:

\begin{equation}\label{eq:der-3}
\begin{aligned}
x^{*} \backslash \mathtt{x}
  &= (x \backslash \mathtt{x}) \cdot x^{*}
  &&\mathsf{D}_{*},\ r = x \\
  &= \varepsilon \cdot x^{*}
  &&\mathsf{D}_{\mathrm{lit}},\ \text{since } \mathtt{x} = x
\end{aligned}
\end{equation}

\emph{Step 2.} The expression is now a concatenation, so the
concatenation rule applies with $r = \varepsilon$ and $s = x^{*}$.
Since $n(\varepsilon) = \mathrm{true}$ by $\mathsf{N}_{\varepsilon}$, it
is the nullable case $\mathsf{D}_{\mathrm{seq}}^{\,n}$ that fires, and that is
what introduces the $+$:

\begin{equation}\label{eq:der-4}
\begin{aligned}
(\varepsilon \cdot x^{*}) \backslash \mathtt{x}
  &= (\varepsilon \backslash \mathtt{x}) \cdot x^{*} \;+\; x^{*} \backslash \mathtt{x}
  &&\mathsf{D}_{\mathrm{seq}}^{\,n},\ r = \varepsilon,\ s = x^{*} \\
  &= \emptyset \cdot x^{*} \;+\; x^{*} \backslash \mathtt{x}
  &&\mathsf{D}_{\varepsilon} \\
  &= \emptyset \cdot x^{*} \;+\; (x \backslash \mathtt{x}) \cdot x^{*}
  &&\mathsf{D}_{*},\ r = x \\
  &= \emptyset \cdot x^{*} \;+\; \varepsilon \cdot x^{*}
  &&\mathsf{D}_{\mathrm{lit}}
\end{aligned}
\end{equation}

\emph{Nullability of the result}, by structural recursion:

\begin{equation}\label{eq:der-5}
\begin{aligned}
n(\emptyset \cdot x^{*} + \varepsilon \cdot x^{*})
&= n(\emptyset \cdot x^{*}) \lor n(\varepsilon \cdot x^{*})
   &&\mathsf{N}_{+},\ r = \emptyset \cdot x^{*},\ s = \varepsilon \cdot x^{*} \\
&= (n(\emptyset) \land n(x^{*})) \lor (n(\varepsilon) \land n(x^{*}))
   &&\mathsf{N}_{\mathrm{seq}},\ \text{twice} \\
&= (\mathrm{false} \land \mathrm{true}) \lor (\mathrm{true} \land \mathrm{true})
   &&\mathsf{N}_{\emptyset},\ \mathsf{N}_{*},\ \mathsf{N}_{\varepsilon} \\
&= \mathrm{true}
\end{aligned}
\end{equation}

so $x^{*}$ matches $\mathtt{xx}$, correctly. This short trace makes two
things visible. First, the expression has grown from two nodes to nine
across two steps, purely from the $+$ that $\mathsf{D}_{\mathrm{seq}}^{\,n}$
introduces at each nullable-prefix step; nothing above bounds this
growth, and a third derivative step would nest a further $+$ inside each
existing summand, reaching fourteen nodes. Second, the identities listed
below collapse the result back to $x^{*}$ itself, the expression the
trace started from, which is why applying them after every step keeps
this particular sequence at a fixed size rather than letting it grow.
The collapse takes three of the four identities in turn:

\begin{equation}\label{eq:der-6}
\begin{aligned}
\emptyset \cdot x^{*} + \varepsilon \cdot x^{*}
  &= \emptyset + \varepsilon \cdot x^{*} &&\mathsf{S}_{1},\ r = x^{*} \\
  &= \varepsilon \cdot x^{*}             &&\mathsf{S}_{3},\ r = \varepsilon \cdot x^{*} \\
  &= x^{*}                               &&\mathsf{S}_{2},\ r = x^{*}
\end{aligned}
\end{equation}


\medskip
\noindent Four identities do this collapsing,
Equation~\eqref{eq:simp-identities}, each a direct consequence of the
semantics of Section~\ref{sec:regex-semantics}:

\begin{equation}\label{eq:simp-identities}
\begin{aligned}
\emptyset \cdot r &= \emptyset &&(\mathsf{S}_{1}) &
\emptyset + r &= r &&(\mathsf{S}_{3}) \\
\varepsilon \cdot r &= r &&(\mathsf{S}_{2}) &
r + r &= r &&(\mathsf{S}_{4}) \quad\text{(idempotent alternation)}
\end{aligned}
\end{equation}

\texttt{src/regex/simplify.rs} implements exactly these four rules,
applied to every subexpression; \texttt{match\_deriv}
(\texttt{src/matchers/match\_deriv.rs}) calls it after every derivative
step, which is why neither worked example in this section ever shows an
expression larger than what is written down. Each step in the trace
below is annotated with the identity it uses.

This is not only a matter of efficiency: Antimirov observes that
Brzozowski's original formulation of the concatenation rule, phrased as
multiplication by the nullability value rather than as the explicit
case split used above, produces an infinite, non-collapsing sequence of
derivatives for expressions such as $(x+y)^{*}$ unless simplification is
applied at every step \citep{antimirov1996}. Eager simplification is
therefore what keeps the finiteness property in practice, not merely
what keeps expressions small.

%

\paragraph{Worked example: $r_{01} = (0+1)^{*} \cdot 0 \cdot 1$}
Recall $r_{01}$ from Section~\ref{sec:regex-syntax}, left-associated as
$((0+1)^{*} \cdot 0) \cdot 1$. Derive it against $w = \mathtt{101}$ one
symbol at a time.


Because every rule below is applied to a concatenation, it is worth
fixing the bindings once. $r_{01}$ has three concatenation nodes, and
matching each against the pattern $r \cdot s$ gives:

\begin{center}
\begin{tabular}{lll}
\hline
Node & $r$ & $s$ \\
\hline
$\big((0+1)^{*} \cdot 0\big) \cdot 1$ (the whole of $r_{01}$) & $(0+1)^{*} \cdot 0$ & $1$ \\
$(0+1)^{*} \cdot 0$ & $(0+1)^{*}$ & $0$ \\
inside $\mathsf{D}_{*}$: the star body & $0+1$ & \\
\hline
\end{tabular}
\end{center}

The first row is the one that is easy to get wrong. $r_{01}$ is
left-associated (Section~\ref{sec:regex-syntax}), so its outermost
operator is the concatenation with the trailing literal $1$, and the
rule's $s$ is that single literal rather than $0 \cdot 1$.

\begin{enumerate}
  \item \textbf{Derive by $\mathtt{1}$.} At the outer node,
  $r = (0+1)^{*} \cdot 0$ and $s = 1$. Testing $r$ for nullability,
  $n\big((0+1)^{*} \cdot 0\big) = n\big((0+1)^{*}\big) \land n(0) =
  \mathrm{true} \land \mathrm{false} = \mathrm{false}$ by
  $\mathsf{N}_{\mathrm{seq}}$, $\mathsf{N}_{*}$ and $\mathsf{N}_{\mathrm{lit}}$:
  it does not match the empty word, since it still owes its trailing
  $0$. So the non-nullable case applies and no second summand appears:
  \begin{equation}\label{eq:der-8}
\begin{aligned}
r_{01} \backslash \mathtt{1}
  &= \big( ((0+1)^{*} \cdot 0) \backslash \mathtt{1} \big) \cdot 1
  &&\mathsf{D}_{\mathrm{seq}},\ r = (0+1)^{*} \cdot 0,\ s = 1
\end{aligned}
\end{equation}

  The inner derivative is taken at the middle node, where
  $r = (0+1)^{*}$ and $s = 0$. That $r$ is a star, so working it out
  needs $\mathsf{D}_{*}$ with star body $0+1$:
  \begin{equation}\label{eq:der-9}
\begin{aligned}
(0+1)^{*} \backslash \mathtt{1}
  &= \big( (0+1) \backslash \mathtt{1} \big) \cdot (0+1)^{*}
  &&\mathsf{D}_{*},\ r = 0+1 \\
  &= \big( (0 \backslash \mathtt{1}) + (1 \backslash \mathtt{1}) \big) \cdot (0+1)^{*}
  &&\mathsf{D}_{+},\ r = 0,\ s = 1 \\
  &= (\emptyset + \varepsilon) \cdot (0+1)^{*}
  &&\mathsf{D}_{\mathrm{lit}},\ \text{twice} \\
  &= \varepsilon \cdot (0+1)^{*}
  &&\mathsf{S}_{3},\ r = \varepsilon \\
  &= (0+1)^{*}
  &&\mathsf{S}_{2},\ r = (0+1)^{*}
\end{aligned}
\end{equation}

  Substituting back into the middle node. Here
  $n\big((0+1)^{*}\big) = \mathrm{true}$ by $\mathsf{N}_{*}$, so this
  time it \emph{is} the nullable case, and a second summand does appear:
  \begin{equation}\label{eq:der-10}
\begin{aligned}
((0+1)^{*} \cdot 0) \backslash \mathtt{1}
  &= \big( (0+1)^{*} \backslash \mathtt{1} \big) \cdot 0 \;+\; (0 \backslash \mathtt{1})
  &&\mathsf{D}_{\mathrm{seq}}^{\,n},\ r = (0+1)^{*},\ s = 0 \\
  &= (0+1)^{*} \cdot 0 \;+\; \emptyset
  &&\text{previous block; } \mathsf{D}_{\mathrm{lit}} \\
  &= (0+1)^{*} \cdot 0
  &&\mathsf{S}_{3},\ r = (0+1)^{*} \cdot 0
\end{aligned}
\end{equation}
  Note that $\mathsf{S}_{3}$ is stated as $\emptyset + r = r$ while the
  expression here is $r + \emptyset$; the identity is applied up to the
  commutativity of $+$, which \texttt{simplify} implements as a separate
  match arm for each side.

  So, substituting once more into the outer node,
  \begin{equation}\label{eq:der-24}
r_{01} \backslash \mathtt{1} \;=\; \big( (0+1)^{*} \cdot 0 \big) \cdot 1 \;=\; r_{01}.
\end{equation}
  A $\mathtt{1}$ that is not yet part of the final $\mathtt{01}$ loops
  back to the start, since $(0+1)^{*}$ can absorb it.

  \item \textbf{Derive by $\mathtt{0}$.} The bindings are the same as in
  step~1, since the expression is unchanged. The computation is also the
  same with $0$ and $1$ exchanged, so $(0+1)^{*} \backslash \mathtt{0}$
  again reduces to $(0+1)^{*}$, this time via
  $(\varepsilon + \emptyset)$ rather than $(\emptyset + \varepsilon)$.
  The one substantive difference is at the middle node's second summand:
  \begin{equation}\label{eq:der-11}
\begin{aligned}
((0+1)^{*} \cdot 0) \backslash \mathtt{0}
  &= \big( (0+1)^{*} \backslash \mathtt{0} \big) \cdot 0 \;+\; (0 \backslash \mathtt{0})
  &&\mathsf{D}_{\mathrm{seq}}^{\,n},\ r = (0+1)^{*},\ s = 0 \\
  &= (0+1)^{*} \cdot 0 \;+\; \varepsilon
  &&\mathsf{D}_{\mathrm{lit}},\ \text{since } \mathtt{0} = 0
\end{aligned}
\end{equation}
  Here $\mathsf{D}_{\mathrm{lit}}$ takes its $x = y$ branch and yields
  $\varepsilon$ where step~1 got $\emptyset$, so no simplification
  applies: $\mathsf{S}_{3}$ needs an $\emptyset$ summand and
  $\mathsf{S}_{4}$ needs two identical summands, and neither holds.
  Substituting into the outer node, which is still the non-nullable
  case:
  \begin{equation}\label{eq:der-25}
r_{01} \backslash \mathtt{1} \backslash \mathtt{0} \;=\; \big( (0+1)^{*} \cdot 0 \;+\; \varepsilon \big) \cdot 1.
\end{equation}
  It reads exactly as it should: after \texttt{"10"}, either the
  trailing $\mathtt{01}$ is still owed in full (the $(0+1)^{*} \cdot 0$
  summand), or the just-consumed $\mathtt{0}$ was already the first half
  of it (the $\varepsilon$ summand); both possibilities stay open until
  the next symbol decides between them.

  \item \textbf{Derive by $\mathtt{1}$.} The outer node's binding has
  changed: $r$ is now $(0+1)^{*} \cdot 0 + \varepsilon$ and $s$ is still
  $1$. This $r$ \emph{is} nullable, since
  $n\big((0+1)^{*} \cdot 0 + \varepsilon\big) =
  n\big((0+1)^{*} \cdot 0\big) \lor n(\varepsilon) =
  \mathrm{false} \lor \mathrm{true} = \mathrm{true}$ by $\mathsf{N}_{+}$
  and $\mathsf{N}_{\varepsilon}$, so for the first time the outer node
  takes the nullable case:
  \begin{equation}\label{eq:der-12}
\begin{aligned}
\big( ((0+1)^{*} \cdot 0 + \varepsilon) \cdot 1 \big) \backslash \mathtt{1}
  &= \big( ((0+1)^{*} \cdot 0 + \varepsilon) \backslash \mathtt{1} \big) \cdot 1
     \;+\; (1 \backslash \mathtt{1})
  &&\mathsf{D}_{\mathrm{seq}}^{\,n}
\end{aligned}
\end{equation}
  The inner derivative is now an alternation rather than a
  concatenation, so $\mathsf{D}_{+}$ applies, and its left operand is
  step~1's computation reused verbatim:
  \begin{equation}\label{eq:der-13}
\begin{aligned}
((0+1)^{*} \cdot 0 + \varepsilon) \backslash \mathtt{1}
  &= \big( ((0+1)^{*} \cdot 0) \backslash \mathtt{1} \big) + (\varepsilon \backslash \mathtt{1})
  &&\mathsf{D}_{+},\ r = (0+1)^{*} \cdot 0,\ s = \varepsilon \\
  &= (0+1)^{*} \cdot 0 \;+\; \emptyset
  &&\text{step~1; } \mathsf{D}_{\varepsilon} \\
  &= (0+1)^{*} \cdot 0
  &&\mathsf{S}_{3},\ r = (0+1)^{*} \cdot 0
\end{aligned}
\end{equation}
  Substituting both parts back, and using
  $1 \backslash \mathtt{1} = \varepsilon$ by $\mathsf{D}_{\mathrm{lit}}$:
  \begin{equation}\label{eq:der-26}
r_{01} \backslash \mathtt{1} \backslash \mathtt{0} \backslash \mathtt{1}
  \;=\; \big( (0+1)^{*} \cdot 0 \big) \cdot 1 \;+\; \varepsilon
  \;=\; r_{01} + \varepsilon.
\end{equation}
\end{enumerate}


$r_{01} + \varepsilon$ is nullable, since $n(r_{01} + \varepsilon) =
n(r_{01}) \lor n(\varepsilon) = \mathrm{false} \lor \mathrm{true} =
\mathrm{true}$ by $\mathsf{N}_{+}$ and $\mathsf{N}_{\varepsilon}$, so $\mathtt{101} \in L(r_{01})$, correctly: $\mathtt{101}$
ends in $\mathtt{01}$. The final $\varepsilon$ summand is exactly what
witnesses this: it says the match was already complete the moment the
third symbol was read, which is precisely what makes the whole
expression nullable.

Now for only the word, $w = \mathtt{10}$.
It shares its first two symbols with $\mathtt{101}$,
so deriving $r_{01}$ by $\mathtt{1}$ then by $\mathtt{0}$ produces
the same intermediate expression as step~2 above,
$\big( (0+1)^{*} \cdot 0 + \varepsilon \big) \cdot 1$, but since $\mathtt{10}$
has only two symbols, this \emph{is} the final result for this word,
not an intermediate one. Checking its nullability directly: the expression
is \emph{not} nullable, since neither summand of $(0+1)^{*} \cdot 0 + \varepsilon$
makes the whole concatenation with $1$ nullable: $(0+1)^{*} \cdot 0$ is not nullable,
and $\varepsilon \cdot 1$ still needs the literal $1$. So $\mathtt{10} \notin L(r_{01})$,
correctly: $\mathtt{10}$ ends in $\mathtt{0}$, not $\mathtt{01}$.

\texttt{src/regex/nullable.rs} and \texttt{src/regex/deriv.rs}
implement $n(r)$ and $r \backslash x$ case by case, one \texttt{Regex}
variant at a time. \texttt{match\_deriv}
(\texttt{src/matchers/match\_deriv.rs}) folds \texttt{deriv} over the
input characters, applies \texttt{simplify} after every step exactly as
annotated in the worked example above, and checks \texttt{nullable} on
the result.

The sequence of derivatives $r_0 \xrightarrow{x_1} r_1
\xrightarrow{x_2} \cdots \xrightarrow{x_n} r_n$ can equivalently be read
as tracing a path through an automaton built lazily, on demand: each
$r_i$ is a state, and $\xrightarrow{x}$ is the transition relation. At
every step there is exactly one current expression, playing the role of
exactly one current automaton state, however large that state's internal
syntax tree happens to be. This is a DFA-style view of the construction,
one that Section~\ref{sec:antimirov-partial-derivatives} contrasts
directly with an NFA-style one.

% 

\section{Antimirov Partial Derivatives}
\label{sec:antimirov-partial-derivatives}

The derivative of Section~\ref{sec:brzozowski-derivatives} keeps the
residual language as a single expression, combining every alternative
reached so far with $+$. Antimirov's partial derivatives take a
different route to the same residual language: rather than combining the
alternatives produced along the way, keep them apart as a \emph{set} of
expressions, one per strand of nondeterminism \citep{antimirov1996}.
Writing $pd(r,x)$ for this set, Equation~\eqref{eq:pd} defines it:

\begin{equation}\label{eq:pd}
\begin{aligned}
pd(\emptyset, x) &= \{\} &&(\mathsf{P}_{\emptyset}) \\
pd(\varepsilon, x) &= \{\} &&(\mathsf{P}_{\varepsilon}) \\
pd(y, x) &=
  \begin{cases} \{\varepsilon\} & x = y \\ \{\} & x \neq y \end{cases}
  &&(\mathsf{P}_{\mathrm{lit}}) \\
pd(r + s, x) &= pd(r,x) \cup pd(s,x) &&(\mathsf{P}_{+}) \\
pd(r^{*}, x) &= \{\, \mathit{smartC}(r', r^{*}) \mid r' \in pd(r,x) \,\}
  &&(\mathsf{P}_{*}) \\
pd(r \cdot s, x) &=
  \begin{cases}
    \{\, \mathit{smartC}(r', s) \mid r' \in pd(r,x) \,\} \cup pd(s,x) & n(r) = \mathrm{true} \\
    \{\, \mathit{smartC}(r', s) \mid r' \in pd(r,x) \,\} & n(r) = \mathrm{false}
  \end{cases}
  &&\begin{array}{l}(\mathsf{P}_{\mathrm{seq}}^{\,n}) \\ (\mathsf{P}_{\mathrm{seq}})\end{array}
\end{aligned}
\end{equation}

where the smart constructor, named $\mathsf{SC}$ when cited below, is
\begin{equation}\label{eq:smartc}
\mathit{smartC}(r', s) =
  \begin{cases}
    s & r' = \varepsilon \\
    r' \cdot s & \text{otherwise}
  \end{cases}
\end{equation}
and avoids introducing a spurious $\varepsilon \cdot s$ whenever a
partial derivative of $r$ is exactly $\varepsilon$. The rule names
mirror Section~\ref{sec:brzozowski-derivatives}'s, so a $\mathsf{P}$
rule and the $\mathsf{D}$ rule with the same subscript are the two
readings of the same case, and the derivations below cite them the same
way, binding $r$ and $s$ explicitly at each application. This keeps every partial derivative a
subexpression of $r$ (up to the concatenations $\mathit{smartC}$
builds), rather than an expression that could grow without bound the way
plain $r \backslash x$ can.


%

\paragraph{Worked example: single-symbol matching.}
Consider $x^{*}$ against $\mathtt{xx}$ again, this time under $pd$. The
base case gives $pd(x,\mathtt{x}) = \{\varepsilon\}$, and the star rule
applies the smart constructor to it:

\begin{equation}\label{eq:der-15}
\begin{aligned}
pd(x^{*},\mathtt{x})
  &= \{\, \mathit{smartC}(r', x^{*}) \mid r' \in pd(x,\mathtt{x}) \,\}
  &&\mathsf{P}_{*},\ \text{star body } r = x \\
  &= \{\, \mathit{smartC}(r', x^{*}) \mid r' \in \{\varepsilon\} \,\}
  &&\mathsf{P}_{\mathrm{lit}},\ \text{since } \mathtt{x} = x \\
  &= \{\mathit{smartC}(\varepsilon, x^{*})\} \\
  &= \{x^{*}\}
  &&\mathsf{SC},\ r' = \varepsilon,\ s = x^{*}
\end{aligned}
\end{equation}

Deriving again by $\mathtt{x}$
repeats exactly the same computation, because the set to derive is still
$\{x^{*}\}$: $pd(x^{*},\mathtt{x},\mathtt{x}) = \{x^{*}\}$. This set
contains $x^{*}$, which is nullable, so $x^{*}$ matches $\mathtt{xx}$,
correctly, which is the same conclusion as the derivative trace, reached here in
two steps that never leave $\{x^{*}\}$. Unlike that trace, which passed
through an expression carrying an $\emptyset$ and an $\varepsilon$ neither
one needed, before collapsing back down, $pd$ never produces either kind
of clutter to begin with: $pd(\emptyset,x) = \{\}$ discards empty
alternatives the moment they arise, and $\mathit{smartC}$ discards a
spurious $\varepsilon \cdot s$ the moment it would be built. No separate
simplification pass is needed for this to hold, it follows directly
from how $pd$ and $\mathit{smartC}$ are defined, not from an additional
rewriting step applied afterwards.

%

\paragraph{Worked example: $r_{01} = (0+1)^{*} \cdot 0 \cdot 1$}
Derive $r_{01}$ against $w = \mathtt{101}$ again, this time tracking a
set of residuals instead of one expression.

\begin{enumerate}
  \item \textbf{Derive by $\mathtt{1}$.} The bindings are the same three
  concatenation nodes as in
  Section~\ref{sec:brzozowski-derivatives}'s table, and the work
  proceeds innermost first. Start with the star body, where
  $\mathsf{P}_{+}$ binds $r = 0$ and $s = 1$:
  \begin{equation}\label{eq:der-16}
\begin{aligned}
pd(0{+}1,\ \mathtt{1})
  &= pd(0,\mathtt{1}) \cup pd(1,\mathtt{1})
  &&\mathsf{P}_{+},\ r = 0,\ s = 1 \\
  &= \{\} \cup \{\varepsilon\}
  &&\mathsf{P}_{\mathrm{lit}},\ \text{both branches} \\
  &= \{\varepsilon\}
\end{aligned}
\end{equation}
  Lifting that through the star, whose body is $r = 0+1$:
  \begin{equation}\label{eq:der-17}
\begin{aligned}
pd\big((0+1)^{*},\ \mathtt{1}\big)
  &= \{\, \mathit{smartC}(r',(0+1)^{*}) \mid r' \in \{\varepsilon\} \,\}
  &&\mathsf{P}_{*},\ r = 0+1 \\
  &= \{(0+1)^{*}\}
  &&\mathsf{SC},\ r' = \varepsilon
\end{aligned}
\end{equation}
  At the middle node $r = (0+1)^{*}$ and $s = 0$, and
  $n\big((0+1)^{*}\big) = \mathrm{true}$ by $\mathsf{N}_{*}$, so the
  nullable case applies and a union with $pd(s,x)$ appears:
  \begin{equation}\label{eq:der-18}
\begin{aligned}
pd\big((0+1)^{*}\cdot0,\ \mathtt{1}\big)
  &= \{\, \mathit{smartC}(r',0) \mid r' \in pd\big((0+1)^{*},\mathtt{1}\big) \,\}
     \;\cup\; pd(0,\mathtt{1})
  &&\mathsf{P}_{\mathrm{seq}}^{\,n} \\
  &= \{\mathit{smartC}((0+1)^{*},0)\} \;\cup\; \{\}
  &&\text{above; } \mathsf{P}_{\mathrm{lit}} \\
  &= \{(0+1)^{*}\cdot0\}
  &&\mathsf{SC},\ r' \neq \varepsilon
\end{aligned}
\end{equation}
  Note that $\mathsf{SC}$ took its second branch here, building the
  concatenation, because $r' = (0+1)^{*}$ is not $\varepsilon$.
  Finally the outer node, where $r = (0+1)^{*}\cdot0$ and $s = 1$. That
  $r$ is not nullable, so no union is added:
  \begin{equation}\label{eq:der-19}
\begin{aligned}
pd(r_{01},\ \mathtt{1})
  &= \{\, \mathit{smartC}(r',1) \mid r' \in \{(0+1)^{*}\cdot0\} \,\}
  &&\mathsf{P}_{\mathrm{seq}},\ r = (0+1)^{*}\cdot0,\ s = 1 \\
  &= \{(0+1)^{*}\cdot0\cdot1\} \;=\; \{r_{01}\}
  &&\mathsf{SC}
\end{aligned}
\end{equation}

  A single residual, identical to $r_{01}$ itself: reading a $\mathtt{1}$
  that is not yet part of the final $\mathtt{01}$ loops back to the
  start, exactly as in the derivative trace, but here as a one-element
  set, rather than as one expression equal to $r_{01}$.

  \item \textbf{Derive by $\mathtt{0}$.} The set is still $\{r_{01}\}$,
  so derive $r_{01}$ by $\mathtt{0}$ with the same three bindings. The
  star computation runs as before and again gives
  $pd\big((0+1)^{*},\mathtt{0}\big) = \{(0+1)^{*}\}$. The difference is
  at the middle node's second operand, where $\mathsf{P}_{\mathrm{lit}}$
  now takes its $x = y$ branch:
  \begin{equation}\label{eq:der-20}
\begin{aligned}
pd\big((0+1)^{*}\cdot0,\ \mathtt{0}\big)
  &= \{\mathit{smartC}((0+1)^{*},0)\} \;\cup\; pd(0,\mathtt{0})
  &&\mathsf{P}_{\mathrm{seq}}^{\,n},\ r = (0+1)^{*},\ s = 0 \\
  &= \{(0+1)^{*}\cdot0\} \;\cup\; \{\varepsilon\}
  &&\mathsf{P}_{\mathrm{lit}},\ \text{since } \mathtt{0} = 0 \\
  &= \{(0+1)^{*}\cdot0,\ \varepsilon\}
\end{aligned}
\end{equation}
  Write $P = \{(0+1)^{*}\cdot0,\ \varepsilon\}$ for the set just
  computed. Substituting it into the outer node, $\mathsf{SC}$ now takes
  a different branch for each of its two elements:
  \begin{equation}\label{eq:der-21}
\begin{aligned}
pd(r_{01},\mathtt{1},\mathtt{0})
  &= \{\, \mathit{smartC}(r',1) \mid r' \in P \,\}
  &&\mathsf{P}_{\mathrm{seq}},\ s = 1 \\
  &= \{(0+1)^{*}\cdot0\cdot1,\ \ 1\}
  &&\mathsf{SC}\text{: second branch, then first} \\
  &= \{r_{01},\ 1\}
\end{aligned}
\end{equation}
  The second element is the literal $1$ rather than $\varepsilon \cdot 1$
  precisely because $\mathsf{SC}$ recognized $r' = \varepsilon$ and
  returned $s$ untouched.

  This is the same information as $\big((0+1)^{*}\cdot0 + \varepsilon\big)\cdot1$
  from the derivative trace. Both say "either the trailing $\mathtt{01}$
  is still owed in full, or the just-consumed $\mathtt{0}$ was already its
  first half", but where the derivative folds the two possibilities
  into one expression with an internal $+$, $pd$ keeps them apart as two
  separate residuals, $r_{01}$ and the literal $1$, with no $+$ at all.
  This is the concrete sense in which $pd$ is NFA-shaped: the two-element
  set is exactly the state set an NFA for $r_{01}$ would be in after
  reading $\mathtt{10}$, one state per strand, rather than one DFA state
  encoding both strands internally.

  \item \textbf{Derive by $\mathtt{1}$.} Each element of $\{r_{01},\ 1\}$
  is derived separately and the results unioned. The first repeats
  step~1 exactly; the second is a bare literal, so only
  $\mathsf{P}_{\mathrm{lit}}$ is needed:
  \begin{equation}\label{eq:der-22}
\begin{aligned}
pd(r_{01},\mathtt{1},\mathtt{0},\mathtt{1})
  &= pd(r_{01},\mathtt{1}) \;\cup\; pd(1,\mathtt{1})
  &&\text{union over the frontier} \\
  &= \{r_{01}\} \;\cup\; \{\varepsilon\}
  &&\text{step~1; } \mathsf{P}_{\mathrm{lit}} \\
  &= \{r_{01},\ \varepsilon\}
\end{aligned}
\end{equation}
\end{enumerate}

%

This set contains $\varepsilon$, so it contains a nullable residual, and
$\mathtt{101} \in L(r_{01})$ the same conclusion as the derivative
trace, reached here by checking whether any residual currently held is
nullable, rather than by checking nullability of one combined expression.

Now for $w = \mathtt{10}$, the same way the derivative trace did: it
shares its first two symbols with $\mathtt{101}$, so deriving by
$\mathtt{1}$ then $\mathtt{0}$ (step~2 above) reaches the same set,
$\{r_{01},\ 1\}$ and since $\mathtt{10}$ has only two symbols, this
\emph{is} the final set for this word. Neither element is nullable:
$n(r_{01}) = \mathrm{false}$, and the literal $1$ is nullable only once
it has been derived into $\varepsilon$, not as itself. So
$\mathtt{10} \notin L(r_{01})$, correctly.

The set of all partial derivatives reachable from a fixed $r$ by
derivation with respect to arbitrary words is always finite, bounded by
the number of literal-symbol occurrences in $r$ \citep{antimirov1996}.
$r_{01}$ has four such occurrences, the two inside $(0+1)^{*}$, and one
each for the trailing $0$ and $1$, so at most four distinct partial
derivatives (plus $r_{01}$ itself) can ever appear, no matter how long a
word $r_{01}$ is derived with respect to. The worked example above
already shows three of those residuals directly: every set encountered,
$\{r_{01}\}$, $\{r_{01}, 1\}$, and $\{r_{01}, \varepsilon\}$, is built
entirely from $r_{01}$, the literal $1$, and $\varepsilon$. Ordinary
derivatives are also finite up to language equivalence. This is
\citet{brzozowski1964}'s own finiteness result, the one that makes the
DFA reading at the end of Section~\ref{sec:brzozowski-derivatives}
possible in the first place, but only once the simplification
identities of Section~\ref{sec:brzozowski-derivatives} are applied to
identify equivalent expressions; the partial-derivative bound holds
already at the syntactic level, with no such identification
needed and the same point the single-symbol example above made concretely
for $x^{*}$. The two bounds also differ in kind, not just in what they
need to hold: Brzozowski's own bound on the number of derivative types
is exponential in the structure of $r$
\citep[Theorem~4.3]{brzozowski1964}, while Antimirov's $\|r\|+1$
bound on partial derivatives is linear in the number of literal
occurrences \citep[Theorem~3.4, p.~309]{antimirov1996}.

$pd$ is implemented inside \texttt{match\_pderiv}
(\texttt{src/matchers/match\_pderiv.rs}), which mirrors the worked
example above directly: rather than tracking one expression, it
maintains a \texttt{HashSet<Regex>}
\citep{rust-hashset,rust-hash} of currently-live residuals, replacing
that set at each character with the union of $pd(r',x)$ over every
residual $r'$ currently in it, and accepting once the input is exhausted
if any surviving residual is nullable, which is exactly the frontier
$\{r_{01}\} \to \{r_{01}, 1\} \to \{r_{01}, \varepsilon\}$ traced by hand
above. This is subset simulation of an NFA whose states are
subexpressions of $r_{01}$, run directly on the syntax rather than on a
separately constructed automaton.

%

\section{Matcher Selection}
\label{sec:matcher-selection}

\texttt{match\_naive}, \texttt{match\_deriv}, and \texttt{match\_pderiv}
share the signature \texttt{fn(\&str, \&Regex) -> bool} despite
implementing three structurally different algorithms: exhaustive search
over substring splits, a single evolving DFA-style expression, and an
evolving NFA-style set of expressions. This uniform interface makes the
three matchers directly interchangeable: the same regular expression and
input can be run through any of them, and any disagreement between their
answers points to a bug in one of the three independently-written
implementations rather than to an ambiguity in what the correct answer is
supposed to be. The same pattern of comparing independently-derived
implementations is reused throughout this thesis for parsers rather than
matchers, wherever multiple constructions are expected to agree.

\texttt{MatcherType} (\texttt{src/matchers/selection.rs}) is an enum with
one variant per matcher and a \texttt{matcher()} method returning the
corresponding function pointer. Its purpose is purely to give the
\texttt{regex-engine} command-line tool (\texttt{src/cli/matcher.rs}) a
single place that maps a \texttt{-{}-matcher} flag value
(\texttt{naive}, \texttt{deriv}, \texttt{pderiv}, or \texttt{all},
defaulting to \texttt{deriv}) on the \texttt{match} subcommand to the
matcher it names, rather than every call site duplicating that mapping
itself. All three matchers can be compared side by side on an arbitrary
input through the CLI directly, which is how the two worked examples of
this chapter can be checked end to end rather than only by hand. Running
each of $x^{*}$ against $\mathtt{xx}$, and $r_{01}$ against $\mathtt{101}$
and $\mathtt{10}$, with \texttt{-{}-matcher all} gives:

\begin{center}
\begin{tabular}{llccc}
\hline
Example & Word & \texttt{naive} & \texttt{deriv} & \texttt{pderiv} \\
\hline
$x^{*}$   & \texttt{xx}  & true  & true  & true \\
$r_{01}$  & \texttt{101} & true  & true  & true \\
$r_{01}$  & \texttt{10}  & false & false & false \\
\hline
\end{tabular}
\end{center}

\begin{verbatim}
cargo run -- match "(0|1)*01" "101" --matcher all
\end{verbatim}
prints \texttt{true} for all three matchers and reproduces the second
row directly. This agrees with the hand-derived conclusions of
Sections~\ref{sec:brzozowski-derivatives}
and~\ref{sec:antimirov-partial-derivatives} in every cell: $x^{*}$
accepts $\mathtt{xx}$, $r_{01}$ accepts $\mathtt{101}$, and $r_{01}$
rejects $\mathtt{10}$ and all three independently-written matchers
reach each conclusion the same way. The same
table-plus-CLI-invocation pattern is used again in
Chapters~\ref{sec:derivative-parser}
and~\ref{sec:partial-derivative-parser}
to check the parsers' worked examples against the \texttt{parse}
subcommand instead of \texttt{match}.

Running these through the CLI does, however, lean on more than the three
matchers themselves. \texttt{match\_naive}, \texttt{match\_deriv}, and
\texttt{match\_pderiv} take only an already-built \texttt{Regex} value and
their shared signature is \texttt{fn(\&str, \&Regex) -> bool}, with no
path from a string like \texttt{"(0|1)*01"} to a \texttt{Regex} anywhere
in them. That translation is entirely the frontend's purpose: the \texttt{match}
subcommand (\texttt{src/cli/matcher.rs}) calls
\texttt{frontend::parse\_pattern} \emph{before} it calls whichever matcher
\texttt{-{}-matcher} selected, turning the concrete syntax into the
\texttt{Regex} AST of Section~\ref{sec:regex-rust-implementation} first.
Chapter~\ref{sec:frontend} covers that parsing and translation in detail.

The POSIX parsers built in Chapters~\ref{sec:derivative-parser}
and~\ref{sec:partial-derivative-parser} follow the same pattern one level
up: rather than a Boolean, each parser returns a parse tree, but all of
them remain selectable through a single uniform interface and are
expected to agree with one another wherever they overlap.